\section{Proofs}\label{app:proofs}

Numbers without the prefix S refer to the article.

\begin{proof}[Proof of Proposition~\ref{prop:stability}]
Subtract the two forward maps and use
$1/\widehat q-1/q=\rho e_s/(q\widehat q)$ to obtain
\eqref{eq:forward-error}. Since $u=\rho D$,
$\widehat D=D+e_t+u e_s-\widehat s e_a$ and
$\widehat\rho-\rho=(u-e_a-\rho\widehat D)/\widehat D$, which gives
\eqref{eq:inverse-error}. The triangle inequality yields
$|\widehat D-D|\le B$ and then \eqref{eq:inverse-bound}.
\end{proof}

\begin{proof}[Proof of Proposition~\ref{prop:identity}]
Subtract $h+S(f-h)$ from $f$ to obtain \eqref{eq:residual-exact}. On the design,
$I-K(K+\lambda I)^{-1}=\lambda(K+\lambda I)^{-1}$, which proves
\eqref{eq:design-residual} and its spectral form. Apply the assumed uniform inequality to
$g=f-h$ for \eqref{eq:inherit}.
\end{proof}

\begin{proof}[Proof of Proposition~\ref{prop:alignment}]
For a common operator the corrected error is $a-b$, so expansion of its squared norm gives
\eqref{eq:alignment}. The error of $p_t$ is $a-tb$; minimize the resulting quadratic on
$[0,1]$ to obtain \eqref{eq:shrink-mean}. For distinct operators,
$f-h-S_1(f-h)=(I-S_0)f-[(S_1-S_0)f+(I-S_1)h]$, after which the same expansion applies.
\end{proof}

\begin{proof}[Proof of Proposition~\ref{prop:sensitivity}]
(i) is the chain rule applied to $\widehat y=S_\lambda y+(I-S_\lambda)h_X(y)$, and the two examples are direct
substitutions. (ii) With the mean and kernel frozen the map is affine with linear part $S_\lambda$; the trace is the
sum of the eigenvalues $\mu_j/(\mu_j+\lambda)$, and the operator norm of a symmetric positive semidefinite matrix is
its largest eigenvalue. (iii) $\widehat y-f_X=-(I-S_\lambda)r+S_\lambda\varepsilon$; the cross term has zero mean,
and both matrices are diagonal in the eigenbasis of $K$. (iv) The correction at $x$ is $a_\lambda(x)^{\top}(y-h_X)$,
linear in $\varepsilon$ with coefficient vector $a_\lambda(x)$.
\end{proof}

\begin{proof}[Proof of Proposition~\ref{prop:selection}]
Let $\widehat{\mathcal R}_j$ and $\mathcal R_j$ be the empirical and population risks.
Hoeffding's inequality \cite{hoeffding} and a union bound give
$|\widehat{\mathcal R}_j-\mathcal R_j|\le
B\sqrt{\log(2M/\delta)/(2m)}$ simultaneously for all candidates with probability at least
$1-\delta$. Apply this once to the empirical minimizer and once to a population minimizer,
and use empirical optimality between the two inequalities.
\end{proof}

\begin{proof}[Proof of Theorem~\ref{thm:transfer}]
Write $q=1-\rho s$, $u=L-a=\rho t/q$, $D=t/q$, $\widehat D=\widehat t+\widehat s(L-\widehat a)$ and $z=\bar\rho$.
Suppose first $\widehat D>0$. The map $F(v)=\widehat Dv-(L-\widehat a)$ is increasing and affine, $z$ is the
projection of its root onto $[0,R]$ and $\rho\in[0,R]$, so $(z-\rho)F(z)\le0$: at an interior root $F(z)=0$; at
$z=R$ the root exceeds $R$, so $F(R)<0$ and $z-\rho\ge0$; at $z=0$ the root is negative, so $F(0)>0$ and
$z-\rho\le0$. Since $L-\widehat a=u-e_a$ and, by \eqref{eq:inverse-error},
$\widehat D=D+e_t+ue_s-\widehat se_a$ with $u=\rho D$,
\[
F(z)=D(z-\rho)+(1-z\widehat s)e_a+ze_t+zue_s ,
\]
hence $D(z-\rho)^2\le|z-\rho|\bigl((1-z\widehat s)|e_a|+z|e_t|+zu|e_s|\bigr)$ and, with $D=t/q$,
\[
|z-\rho|\le q(1-z\widehat s)\frac{|e_a|}{t}+qz\frac{|e_t|}{t}+z\rho\,|e_s| .
\]
Now $q(1-z\widehat s)$ is positive and at most $1$ under the hypotheses, which is all that is used here,
$qz\le R$ and $z\rho\le R^2$, so $|z-\rho|\le e_R/t$; and $|z-\rho|\le R$
because both lie in $[0,R]$. If $\widehat D\le0$, then $\widehat t\le\widehat s(e_a-u)\le\widehat s|e_a|$ and
$t=\widehat t-e_t\le|e_t|+\widehat s|e_a|$, so $Rt\le R|e_t|+R\widehat s|e_a|\le R|e_t|+|e_a|\le e_R$; the fallback
value and $\rho$ lie in $[0,R]$, so $|z-\rho|\le R=\min\{R,e_R/t\}$. For optimality take $s=\widehat s=0$,
$\widehat t=t$, $\rho=r_0$ and $\widehat a=a+ct$ with $0<c<r_0<R$: then $|\bar\rho-\rho|=c=e_R/t$ exactly, with a
positive denominator and no projection active. This proves (i). For (ii) integrate (i), split at $t=\tau$ and bound
$e_R^2/t^2$ by $e_R^2/\tau^2$ on $\{t>\tau\}$. For (iii), $\varepsilon_R=0$ gives zero error by (i). Otherwise take $\tau=\varepsilon_R^{2/(\beta+2)}\le u_0$; then
$R^2A\tau^\beta+\varepsilon_R^2/\tau^2=(R^2A+1)\varepsilon_R^{2\beta/(\beta+2)}$. (iv) is (i) with $t\ge t_0$.
\end{proof}

\begin{proof}[Proof of Proposition~\ref{prop:attained}]
$L=r_0t$ and $\widehat\rho=(r_0t-\widehat a)/t$; $e_R=|\widehat a|=ct\,\1\{t\le\tau\}$; with density
$\beta u^{\beta-1}$ on $(0,1)$, $\E_P[t^2\1\{t\le\tau\}]=\beta\tau^{\beta+2}/(\beta+2)$; the retrieval error equals
$c$ exactly on $\{t\le\tau\}$, an event of probability $\tau^\beta$. Eliminate $\tau$.
\end{proof}

\begin{proof}[Proof of Theorem~\ref{thm:profile}]
If $R=0$ the statement is immediate; assume $R>0$. Conditional on $t$ the squared error is at most both $R^2$ and $e_R^2/t^2$ by Theorem~\ref{thm:transfer}(i), so
$\E[|\bar\rho-\rho|^2\mid t]\le\min\{R^2,\E[e_R^2\mid t]/t^2\}\le\min\{R^2,\sigma^2t^{2\nu-2}\}$; integrate. With
$t_*=(\sigma/R)^{1/(1-\nu)}$,
$\E_P\min\{R^2,\sigma^2t^{2\nu-2}\}=R^2F_t(t_*)+\sigma^2\E_P[t^{2\nu-2}\1\{t>t_*\}]$; the first term is at most
$R^2At_*^\beta=O(\sigma^{\beta/(1-\nu)})$ once $t_*\le u_0$, and
$\E_P[t^{2\nu-2}\1\{t>t_*\}]=\int_0^{t_*^{2\nu-2}}\Prob(t_*<t<v^{-1/(2-2\nu)})\,dv
\le u_0^{2\nu-2}+A\int_{u_0^{2\nu-2}}^{t_*^{2\nu-2}}v^{-\beta/(2-2\nu)}\,dv$, which is $O(1)$ for
$\beta>2(1-\nu)$, $O(\log(1/\sigma))$ for $\beta=2(1-\nu)$ and $O(t_*^{\beta-2(1-\nu)})$ for $\beta<2(1-\nu)$. At
$\nu=1$ the bound is $\sigma^2$ before any tail enters. In the model,
$c^{-2}\E_P|\bar\rho-\rho|^2=\E_P\min\{1,(\sigma/c)^2t^{2\nu-2}\}$, the same integral with equality.
\end{proof}

\begin{proof}[Proof of Proposition~\ref{prop:stackweights}]
(i) With $a=s=0$ and exact $t$ and $s$, the stack's path-radiance error is
$e(w)=((1-w)b,\,wc\delta)$, and the radiance at $\rho$ is predicted with the same error, since
$q=\widehat q=1$; this gives $F$. The constrained retrieval evaluates
$(L-\widehat a)/t=\rho-e_i/t_i$ at entry $i$, with denominator $t_i>0$, and
$|e_i|/t_i\le\max\{b,c\}\le\min\{\rho,R-\rho\}$ keeps it in $[0,R]$, so the projection is inactive
and $\bar\rho-\rho=-((1-w)b,\,wc)$; this gives $E$. Both are convex quadratics in $w$, and setting
their derivatives to zero gives $w_f$ and $w_r$. Substituting,
\[
E(w_f)=\frac{b^2c^2\,(c^2\delta^4+b^2)}{2(b^2+c^2\delta^2)^2},\qquad
E(w_r)=\frac{b^2c^2}{2(b^2+c^2)},
\]
whose ratio is the one stated. It tends to $(b^2+c^2)/b^2$ as $\delta\to0$; at $c=b/\delta$ the
two factors of the numerator are $b^2(1+\delta^2)$ and $b^2(1+\delta^2)/\delta^2$ and the
denominator is $4b^4$. Since $w_f$ minimizes $F$, $F(w_f)\le F(w_r)$. Finally, with
$\widetilde e_t=\widetilde e_s=0$ and $\tau\le\delta\le1$, the weights $(t_i\vee\tau)^{-2}$ in
\eqref{eq:Jtau} are $1$ and $\delta^{-2}$, so $J_\tau(w)=\tfrac12[(1-w)^2b^2+w^2c^2\delta^2/\delta^2]=E(w)$.

(ii) With $s=\widehat s=0$ and $\widehat t=t$, the predicted radiance at $\rho$ is
$\widehat a+\rho t$, so its error is the path-radiance error $r_w$ of the combination, and the
retrieval evaluates $(L-\widehat a)/t=\rho-r_w/t$, which lies in $[0,R]$ by hypothesis. Hence
$F(w)=\E_Pr_w^2$ and $E(w)=\E_P[r_w^2/t^2]$. Since the weights sum to one, $r_w$ lies in
$\mathcal L$, so $\lambda_{\min}F(w)\le E(w)\le\lambda_{\max}F(w)$ for every $w$ (both sides vanish
when $\E_Pr_w^2=0$), and
\[
E(w_f)\le\lambda_{\max}F(w_f)\le\lambda_{\max}F(w_r)\le\frac{\lambda_{\max}}{\lambda_{\min}}\,E(w_r).
\]
If $t_{\min}\le t\le t_{\max}$ almost surely, then
$r^2/t_{\max}^2\le r^2/t^2\le r^2/t_{\min}^2$ pointwise, so $\lambda_{\max}\le t_{\min}^{-2}$ and
$\lambda_{\min}\ge t_{\max}^{-2}$. In the example of (i), $\mathcal L=\mathbb R^2$ and the ratio
$\tfrac12(r_1^2+r_2^2/\delta^2)/\tfrac12(r_1^2+r_2^2)$ runs from $1$ to $\delta^{-2}$, while
$E(w_f)/E(w_r)$ at $c=b/\delta$ is at least $\delta^{-2}/4$.
\end{proof}

\begin{proof}[Proof of Proposition~\ref{prop:perturb}]
$\widetilde H=H^{1/2}(I+B)H^{1/2}$ gives $\widetilde H^{-1}=H^{-1/2}QH^{-1/2}$ and $Q=I-QB$. Hence
$\widetilde k_x^{\top}\widetilde H^{-1}(Y-\widetilde M)-k_x^{\top}H^{-1}(Y-M)
=\widetilde g^{\top}Q(W-U)-g^{\top}W
=\widetilde g^{\top}W-\widetilde g^{\top}QBW-\widetilde g^{\top}QU-g^{\top}W
=v^{\top}W-(Q\widetilde g)^{\top}(BW+U)$, which is (i) since $Q$ is symmetric; (ii) bounds its two terms. For (iii)
write the same difference as $v^{\top}QW+g^{\top}(Q-I)W-\widetilde g^{\top}QU$ and use
$\|Q\|_{\mathrm{op}}=\gamma_Q$ and $\|Q-I\|_{\mathrm{op}}=\|QB\|_{\mathrm{op}}=\eta_Q$. For (iv), $v=0$, $U=0$ and
$B=sH^{-1}$ has eigenvalues $s/\lambda_i(H)\ge0$, so $\eta_Q=\max_i(s/\lambda_i)/(1+s/\lambda_i)=s/(\lambda_{\min}(H)+s)$.
\end{proof}

